\documentclass[letterpaper,fleqn]{article}
\usepackage[spanish,es-noshorthands]{babel}
\usepackage[utf8]{inputenc} 
\usepackage[left=1cm, right=1cm, top=1.5cm, bottom=1.7cm]{geometry}
\usepackage{mathexam}
\usepackage{amsmath}
\usepackage{graphicx}

\ExamClass{\includegraphics[height=16pt]{Images/logo-sed.png} Aritmética $6^{\circ}$}
\ExamName{Sustentación Recomendaciones III}
\ExamHead{\includegraphics[height=16pt]{Images/logo-colegio.png} IEDAB}
\newcommand{\LineaNombre}{%
\par
\vspace{\baselineskip}
Nombre:\hrulefill \; Curso: \underline{\hspace*{48pt}} \; Fecha: \underline{\hspace*{2.5cm}} \relax
\par}
\let\ds\displaystyle

\begin{document}
\ExamInstrBox{
Respuesta sin justificar mediante procedimiento no será tenida en cuenta en la calificación. Escriba sus respuestas en el espacio indicado. Tiene 45 minutos para contestar esta prueba.}
\LineaNombre
\begin{enumerate}
 \item Dados los números 96, 112 y 104:
 \begin{enumerate}
 \item Escríbalos como producto de factores primos\noanswer
 \item ¿Cuáles son sus factores comunes?\noanswer
 \item ¿Cuál es el máximo común divisor y el mínimo común múltiplo de ellos.\noanswer
 \end{enumerate}
 \item En un campamento hay 83 niños y niñas. ¿Qué problema tienen para hacer equipos con el mismo número de integrantes?\noanswer
 \item Calcule por el método que prefiera el máximo común divisor y el mínimo común múltiplo de los siguientes grupos de números:
 \begin{enumerate}
 \item 108, 198 y 136\noanswer
  \newpage
 \item 100, 84 y 91\noanswer
 \item 90, 216 y 198\noanswer
 \end{enumerate}
 \item Pablo y Nuria son primos y visitan a sus abuelos a menudo: Pablo va cada 8 días y Nuria cada 10. Si coincidieron en casa de los abuelos el día de Navidad. ¿Cuándo se volverán a encontrar? ¿Cuántas visitas habrá hecho cada uno?\noanswer
\item En una clase hay 24 alumnos y en otra 32. Para hacer una actividad de geografía se forman en cada clase grupos del mismo número de alumnos, de manera que haya el menor número de grupos posibles. ¿Cuántos alumnos componen cada grupo? ¿Cuántos grupos se forman en total?\noanswer
 \end{enumerate}

\end{document}
